%% pc-pile-accumulateur — accumulateur au plomb en décharge (à gauche) et en recharge (à droite).
%% Même géométrie dans les deux panneaux : les polarités ⊕ (électrode PbO2) et ⊖ (électrode Pb)
%% NE CHANGENT PAS ; ce qui s'inverse, c'est le sens du courant i et celui des électrons,
%% donc le rôle anode / cathode de chaque électrode.
%% Bleu : réduction et courant conventionnel ; rouge : oxydation ; vert : électrons et pont salin.
\documentclass[tikz,border=4pt]{standalone}
\usepackage{liaisons}
\begin{document}
\begin{tikzpicture}[x=1cm, y=1cm,
  fl/.style={-{Stealth[length=5pt]}, line width=0.9pt},
  fil/.style={line width=1pt},
  eti/.style={font=\scriptsize, inner sep=2pt, align=center},
  bloc/.style={font=\scriptsize, inner sep=1.5pt, align=center},
  titre/.style={font=\small\bfseries}]

%% ================= éléments communs à un panneau =================================
%% bécher : \becher{x gauche}
\newcommand\becher[1]{%
  \fill[solideB!8] (#1,0.05) rectangle (#1+1.7,1.35);
  \draw[line width=0.8pt] (#1,1.6) -- (#1,0) -- (#1+1.7,0) -- (#1+1.7,1.6);
  \draw[solideB!45, line width=0.6pt] (#1,1.35) -- (#1+1.7,1.35);}
%% électrode : \electrode{x gauche}{formule}
\newcommand\electrode[2]{%
  \fill[black!35] (#1,0.35) rectangle (#1+0.3,2.9);
  \draw[line width=0.5pt, black!60] (#1,0.35) rectangle (#1+0.3,2.9);
  \node[eti, rotate=90, text=white, inner sep=1pt] at (#1+0.15,1.05) {#2};}
%% ossature d'une cellule (béchers, électrodes, pont salin, fils, polarités)
\newcommand\cellule{%
  \becher{0.15} \becher{2.75}
  \electrode{0.60}{$\mathrm{PbO_2}$} \electrode{3.70}{$\mathrm{Pb}$}
  %% pont salin
  \draw[solideE, line width=3.5pt, line join=round]
    (1.45,1.05) -- (1.45,2.10) -- (3.15,2.10) -- (3.15,1.05);
  \node[eti, text=solideE, anchor=south, inner sep=3pt] at (2.30,2.16) {pont salin};
  %% polarités : elles ne changent pas d'un panneau à l'autre
  \fill[solideA] (0.30,2.60) circle (0.19);
  \node[font=\scriptsize\bfseries, text=white] at (0.30,2.60) {$+$};
  \fill[black] (4.30,2.60) circle (0.19);
  \node[font=\scriptsize\bfseries, text=white] at (4.30,2.60) {$-$};
  %% fils du circuit extérieur
  \draw[fil] (0.75,2.90) -- (0.75,3.90) -- (1.45,3.90);
  \draw[fil] (3.85,2.90) -- (3.85,3.90) -- (3.15,3.90);
  \draw[line width=0.5pt, black!60] (0.49,2.60) -- (0.60,2.60);
  \draw[line width=0.5pt, black!60] (4.00,2.60) -- (4.11,2.60);}
%% flèches du circuit : \fleches{1} = courant de gauche à droite (décharge), {-1} sinon
\newcommand\fleches[1]{%
  \ifnum#1=1
    \draw[fl, solideB] (0.80,4.06) -- (1.33,4.06);
    \draw[fl, solideB] (3.27,4.06) -- (3.80,4.06);
    \draw[fl, solideE] (1.33,3.74) -- (0.80,3.74);
    \draw[fl, solideE] (3.80,3.74) -- (3.27,3.74);
  \else
    \draw[fl, solideB] (1.33,4.06) -- (0.80,4.06);
    \draw[fl, solideB] (3.80,4.06) -- (3.27,4.06);
    \draw[fl, solideE] (0.80,3.74) -- (1.33,3.74);
    \draw[fl, solideE] (3.27,3.74) -- (3.80,3.74);
  \fi
  \node[eti, text=solideB, anchor=south, inner sep=1pt] at (1.06,4.12) {$i$};
  \node[eti, text=solideE, anchor=north, inner sep=1pt] at (1.06,3.68) {$e^-$};}

%% ================= panneau gauche : DÉCHARGE =====================================
\begin{scope}[shift={(0,0)}]
  \node[titre] at (2.30,4.75) {Décharge (l'accumulateur débite)};
  \cellule
  \fleches{1}
  %% récepteur
  \draw[fil, fill=white] (1.43,3.62) rectangle (3.17,4.18);
  \node[font=\tiny, inner sep=1pt, align=center] at (2.30,3.90) {charge\\électrique};
  %% demi-équations
  \node[bloc, text=solideB, anchor=north] at (2.30,-0.30)
    {$\mathrm{PbO_2}$ ($\oplus$) — \textbf{cathode}, réduction :\\
     $\mathrm{PbO_2} + 4\,\mathrm{H^+} + 2\,e^- = \mathrm{Pb^{2+}} + 2\,\mathrm{H_2O}$};
  \node[bloc, text=solideA, anchor=north] at (2.30,-1.15)
    {$\mathrm{Pb}$ ($\ominus$) — \textbf{anode}, oxydation :\\
     $\mathrm{Pb} = \mathrm{Pb^{2+}} + 2\,e^-$};
\end{scope}

%% ================= filet de séparation ===========================================
\draw[black!30, line width=0.6pt] (5.00,-1.95) -- (5.00,4.60);

%% ================= panneau droit : RECHARGE ======================================
\begin{scope}[shift={(5.40,0)}]
  \node[titre] at (2.30,4.75) {Recharge (sur un générateur)};
  \cellule
  \fleches{-1}
  %% générateur : sa borne + alimente l'électrode ⊕ de l'accumulateur
  \draw[fil, fill=white] (1.43,3.62) rectangle (3.17,4.18);
  \node[font=\tiny, inner sep=1pt] at (2.30,3.90) {générateur};
  \node[font=\scriptsize\bfseries, text=solideA] at (1.31,4.32) {$+$};
  \node[font=\scriptsize\bfseries] at (3.29,4.32) {$-$};
  %% demi-équations (les rôles anode / cathode sont échangés)
  \node[bloc, text=solideA, anchor=north] at (2.30,-0.30)
    {$\mathrm{PbO_2}$ ($\oplus$) — \textbf{anode}, oxydation :\\
     $\mathrm{Pb^{2+}} + 2\,\mathrm{H_2O} = \mathrm{PbO_2} + 4\,\mathrm{H^+} + 2\,e^-$};
  \node[bloc, text=solideB, anchor=north] at (2.30,-1.15)
    {$\mathrm{Pb}$ ($\ominus$) — \textbf{cathode}, réduction :\\
     $\mathrm{Pb^{2+}} + 2\,e^- = \mathrm{Pb}$};
\end{scope}

%% ================= équation réversible ============================================
\node[font=\small, anchor=east] at (3.85,-2.70)
  {$\mathrm{PbO_2} + 4\,\mathrm{H^+} + \mathrm{Pb}$};
\node[font=\small, anchor=west] at (5.95,-2.70)
  {$2\,\mathrm{Pb^{2+}} + 2\,\mathrm{H_2O}$};
\draw[fl, solideA] (4.00,-2.58) -- (5.80,-2.58);
\node[font=\scriptsize, text=solideA, anchor=south, inner sep=2pt] at (4.90,-2.54) {décharge};
\draw[fl] (5.80,-2.86) -- (4.00,-2.86);
\node[font=\scriptsize, anchor=north, inner sep=2pt] at (4.90,-2.90) {recharge};

%% ================= bandeau ========================================================
\node[draw, line width=0.6pt, rounded corners=2pt, fill=black!3, inner sep=5pt,
      font=\small, align=center] at (4.95,-3.83)
  {Les polarités $\oplus$ et $\ominus$ ne changent pas ; le sens des électrons, si.};
\end{tikzpicture}
\end{document}
